% Source: public/notes/物工/問題/2019ap_exam.pdf, PDF pp. 6--7.
\documentclass[11pt,a4paper]{article}
\usepackage[margin=24mm]{geometry}
\usepackage{fontspec}
\usepackage{xeCJK}
\setCJKmainfont{HaranoAjiMincho-Regular.otf}
\setCJKsansfont{HaranoAjiGothic-Medium.otf}
\usepackage{amsmath,amssymb}
\usepackage{tikz}
\usetikzlibrary{arrows.meta}
\setlength{\parindent}{1em}
\setlength{\parskip}{0.35em}
\renewcommand{\figurename}{図}
\title{物理工学専攻 院試解答集\\2019年 物理学 第3問}
\author{政岡凜太郎}
\date{}
\begin{document}
\maketitle

\section*{[1] 単原子理想気体の分配関数}
古典気体では、$N$ 個の同種粒子の分配関数は位相空間積分を
$N!(2\pi\hbar)^{3N}$ で割ったものである。運動量を
$\boldsymbol{p}_j=(p_{jx},p_{jy},p_{jz})$ とすると、
\begin{align}
 Z^{(g)}(V,\beta,N)
 &=\frac{1}{N!(2\pi\hbar)^{3N}}
   \int_{V^N}\mathrm{d}^{3N}r\int_{\mathbb{R}^{3N}}\mathrm{d}^{3N}p\,
   \exp\!\left[-\beta\sum_{j=1}^{N}\frac{\boldsymbol{p}_j^2}{2m}\right]\\
 &=\frac{V^N}{N!(2\pi\hbar)^{3N}}
   \left(\int_{-\infty}^{\infty}
        \exp\!\left(-\frac{\beta p^2}{2m}\right)\mathrm{d}p\right)^{3N}\\
 &=\frac{V^N}{N!(2\pi\hbar)^{3N}}
   \left(\frac{2\pi m}{\beta}\right)^{3N/2}\\
 &=\frac{V^N}{N!}
   \left(\frac{m}{2\pi\hbar^2\beta}\right)^{3N/2}.
 \label{eq:canonical}
\end{align}
ここで1成分ごとのガウス積分
$\int_{-\infty}^{\infty}e^{-ap^2}\mathrm{d}p=\sqrt{\pi/a}$
を用いた。

\section*{[2] 気体の大分配関数}
\begin{equation}
 z(\beta,\mu)=e^{\beta\mu}
     \left(\frac{m}{2\pi\hbar^2\beta}\right)^{3/2}
\end{equation}
とおけば、式\eqref{eq:canonical}より
\begin{align}
 Z_G^{(g)}(V,\beta,\mu)
 &=\sum_{N=0}^{\infty}\frac{[Vz(\beta,\mu)]^N}{N!}\\
 &=\exp\!\left[V e^{\beta\mu}
        \left(\frac{m}{2\pi\hbar^2\beta}\right)^{3/2}\right].
 \label{eq:grand-gas}
\end{align}
これは指数関数の級数展開そのものである。

\section*{[3] 圧力}
式\eqref{eq:grand-gas}を問題で与えられた式に代入すると
\begin{align}
 P(\beta,\mu)
 &=\frac{1}{\beta}\frac{\partial}{\partial V}
       \log Z_G^{(g)}(V,\beta,\mu)\\
 &=\frac{e^{\beta\mu}}{\beta}
       \left(\frac{m}{2\pi\hbar^2\beta}\right)^{3/2}
 =e^{\beta\mu}\left(\frac{m}{2\pi\hbar^2}\right)^{3/2}\beta^{-5/2}.
 \label{eq:pressure}
\end{align}

\section*{[4] 1サイトの大分配関数}
大分配重みは $e^{-\beta(E-\mu N)}$ である。空サイトは $(E,N)=(0,0)$
なので重みは $1$、吸着状態は $(E,N)=(-\epsilon,1)$ なので重みは
$e^{\beta(\epsilon+\mu)}$ となる。従って
\begin{equation}
 \xi_G^{(a)}=1+e^{\beta(\epsilon+\mu)}.
 \label{eq:site-partition}
\end{equation}

\section*{[5] 吸着原子密度}
サイトは独立なので全体の大分配関数は
$Z_G^{(a)}=(\xi_G^{(a)})^{N_a}$ である。大分配関数の化学ポテンシャル
微分から平均粒子数を求めると
\begin{align}
 n_a
 &=\frac{1}{N_a\beta}\frac{\partial}{\partial\mu}
             \log Z_G^{(a)}\\
 &=\frac{1}{\beta}\frac{\partial}{\partial\mu}
             \log\!\left(1+e^{\beta(\epsilon+\mu)}\right)\\
 &=\frac{e^{\beta(\epsilon+\mu)}}{1+e^{\beta(\epsilon+\mu)}}
 =\frac{1}{1+e^{-\beta(\epsilon+\mu)}}.
 \label{eq:occupancy-mu}
\end{align}
この値は1サイトあたりの平均占有数なので、$0<n_a<1$ である。

\section*{[6] 圧力と温度による表示}
式\eqref{eq:pressure}から
\begin{equation}
 e^{-\beta\mu}
 =\frac{1}{P}\left(\frac{m}{2\pi\hbar^2}\right)^{3/2}
                  \beta^{-5/2}.
\end{equation}
これを式\eqref{eq:occupancy-mu}に代入し、$\beta=1/(k_BT)$ とすれば
\begin{align}
 n_a(P,T)
 &=\frac{1}{1+\displaystyle\frac{1}{P}
       \left(\frac{m}{2\pi\hbar^2}\right)^{3/2}
       (k_BT)^{5/2}e^{-\epsilon/(k_BT)}}\\
 &=\frac{P}{P+P_*(T)},\\
 P_*(T)&=\left(\frac{m}{2\pi\hbar^2}\right)^{3/2}
       (k_BT)^{5/2}e^{-\epsilon/(k_BT)}.
 \label{eq:occupancy-PT}
\end{align}
$P_*(T)$ は、密度を $1/2$ にする圧力の尺度である。

\section*{[7] 圧力・温度依存性}
温度一定では $P_*(T)$ は定数だから、$n_a=P/[P+P_*(T)]$ は圧力とともに
単調増加する。$P\to0$ で0、$P\to\infty$ で1に近づく。

温度依存性は $P_*(T)$ の変化で決まる。その対数微分は
\begin{equation}
 \frac{\mathrm{d}}{\mathrm{d}T}\log P_*(T)
 =\frac{5}{2T}+\frac{\epsilon}{k_BT^2}>0,
 \qquad T>0.
\end{equation}
従って一定圧力では $P_*(T)$ が厳密に増加し、$n_a$ は厳密に減少する。
$T\to0^+$ では $P_*(T)\to0$ なので $n_a\to1$、$T\to\infty$ では
$P_*(T)\to\infty$ なので $n_a\to0$ である。図では
$T_*=\epsilon/k_B$ および $P_{\mathrm{ref}}=P_*(T_*)$ を基準に取る。
式\eqref{eq:occupancy-PT}より
$P_{\mathrm{ref}}=(m/(2\pi\hbar^2))^{3/2}\epsilon^{5/2}e^{-1}$ である。
左図は $T=T_*$、右図は $P=P_{\mathrm{ref}}$ の例である。

\begin{figure}[htbp]
\centering
\begin{tikzpicture}[x=1.05cm,y=3.0cm,>=Stealth,font=\small]
  \draw[->] (0,0)--(5.1,0) node[right] {$P/P_{\mathrm{ref}}$};
  \draw[->] (0,0)--(0,1.14) node[above] {$n_a$};
  \draw[dashed,gray] (0,1)--(4.8,1);
  \draw[thick,blue,domain=0:4.8,samples=100,variable=\p]
    plot (\p,{\p/(1+\p)});
  \draw (1,0.02)--(1,-0.02) node[below=2pt] {$1$};
  \draw (0.04,0.5)--(-0.04,0.5) node[left=2pt] {$1/2$};
  \node at (2.8,0.29) {$T=T_*$};
\end{tikzpicture}
\hspace{1.2cm}
\begin{tikzpicture}[x=1.05cm,y=3.0cm,>=Stealth,font=\small]
  \draw[->] (0,0)--(5.1,0) node[right] {$T/T_*$};
  \draw[->] (0,0)--(0,1.14) node[above] {$n_a$};
  \draw[dashed,gray] (0,1)--(4.8,1);
  \draw[thick,red,domain=0.06:4.8,samples=140,variable=\t]
    plot (\t,{1/(1+(\t^2.5)*exp(1-1/\t))});
  \draw (1,0.02)--(1,-0.02) node[below=2pt] {$1$};
  \draw (0.04,0.5)--(-0.04,0.5) node[left=2pt] {$1/2$};
  \node at (2.65,0.72) {$P=P_{\mathrm{ref}}$};
\end{tikzpicture}
\caption{吸着密度の圧力・温度依存性。左では $T=T_*$、右では $P=P_{\mathrm{ref}}$ を固定した。}
\label{fig:adsorption}
\end{figure}
\end{document}
